\chapter{绪论}

\section{研究背景与意义}

考勤管理是各类企事业单位、学校等组织日常运营中的基础性工作。长期以来，人员考勤主要依赖人工签到、纸质登记、IC卡刷卡和指纹采集等手段。人工签到和纸质登记方式存在效率低下、数据易被篡改、统计工作繁琐等问题，面对较大规模人员群体时，排队等候时间显著增加，给使用者带来不便。IC卡考勤虽然在一定程度上提升了签到效率，但卡片遗失、损坏、遗忘等现象频发，且代刷卡行为难以有效防范，管理人员事后核查成本较高。指纹考勤需要接触式采集，在公共卫生意识日益提高的背景下，多人共用同一采集设备存在卫生隐患，同时指纹传感器对采集角度和手指状态较为敏感，在干燥、潮湿或磨损情况下识别率明显下降。

人脸识别技术因其非接触、自然交互、难以代刷等特点，逐渐成为考勤领域的研究热点和应用方向。基于摄像头采集的人脸图像，系统可以在用户无感知的情况下完成身份确认，无需携带额外介质，也不强制要求特定的身体接触。在技术成熟度方面，近年来深度学习驱动的人脸检测和特征提取方法取得了显著进展，在公开数据集上的识别准确率已达到实用级别。与此同时，边缘计算设备的算力也在不断提升，使得在低成本嵌入式设备上运行人脸检测和识别模型成为可能。

树莓派作为一款价格低廉、体积小巧的单板计算机，具备完整的GPIO、USB、网络等接口，运行Linux系统，能够部署Python生态中的深度学习推理框架。将人脸识别考勤系统部署在树莓派上，具有以下几方面优势：其一，部署成本低，一台树莓派加上USB摄像头即可构成完整的考勤终端，不需要专门的工控机或服务器；其二，本地化部署，数据存储在设备端，不依赖外部云服务，有利于隐私保护；其三，功耗极低，适合长期不间断运行；其四，系统可独立工作，不依赖外部网络，在网络条件不稳定的场所也能正常使用。

本文的研究目标是设计并实现一套部署在树莓派端的人脸识别考勤系统，打通从成员名单导入、用户注册、人脸采集、管理员审核、实时识别到自动签到的完整业务流程。该系统面向真实考勤场景，需要在有限算力下保证识别准确性和响应速度，同时需要防范照片、屏幕翻拍、预录视频等常见攻击手段。系统的实现不仅是对人脸识别技术在考勤场景中的一次完整工程实践，也为低成本边缘设备上运行AI推理任务提供了可参考的方案。

\section{国内外研究现状}

在人脸识别技术方面，Parkhi等\cite{parkhi2015deep}提出了VGG-Face模型，通过在大规模人脸数据集上进行预训练，证明了深度卷积神经网络在人脸识别任务中的有效性。Schroff等\cite{schroff2015facenet}提出的FaceNet模型首次将三元组损失（Triplet Loss）引入人脸识别，直接优化特征向量间的欧氏距离，为后续研究奠定了基础。Deng等\cite{deng2019arcface}提出的ArcFace通过在特征空间中引入加性角度间隔损失，进一步提升了模型的判别能力。Wang等\cite{wang2018cosface}提出的CosFace则通过大间隔余弦损失实现了类似的效果。这些方法构成了当前人脸识别领域的技术基础。

在人脸检测方面，传统方法如Viola-Jones基于Haar特征的级联分类器在正面光照充足条件下表现尚可，但在侧脸、遮挡和低照度场景下性能急剧下降。Zhang等\cite{zhang2016mtcnn}提出的MTCNN通过三阶段级联网络实现人脸检测和关键点定位，成为业界广泛采用的方案之一。然而，MTCNN的三阶段推理架构在算力受限的设备上难以满足实时性要求。OpenCV官方推出的YuNet模型\cite{yunet}针对端侧部署场景进行了轻量化设计，采用单阶段密集检测思路，模型体积小、推理速度快，适合树莓派等算力受限的设备运行。

活体检测技术是防止人脸系统被照片、视频等伪造手段欺骗的关键环节。当前活体检测方法可分为主动式和被动式两类。主动式方法要求用户配合完成特定动作（如眨眼、转头），通过分析动作执行过程中的面部变化判断活体，安全性较高但用户体验较差。被动式方法无需用户配合，通过分析图像的纹理特征（如摩尔纹、屏幕反光）或生物信号（如rPPG远程光电容积脉搏波）判断活体，用户体验更好但技术难度更大。Zhang等\cite{silentface}提出的Silent-Face-Anti-Spoofing模型基于MobileNetV3架构，可在端侧CPU上实现实时推理。Liu等\cite{li2025survey}的综述指出，多帧时序信息融合是被动活体检测的重要趋势。George和Marcel\cite{george2019pad}提出了基于像素级二值监督的人脸呈现攻击检测方法。近年来的研究进一步探索了基于rPPG生理信号的活体检测\cite{pmc2022rppg,acm2014rppg}，以及基于自拍动态分析的被动活体检测\cite{arxiv2025selfie}。

在边缘AI部署方面，将深度学习模型部署到资源受限设备是一个活跃的研究方向。Zhang等\cite{acm2025lightweight}在树莓派环境下设计了轻量级人脸识别考勤系统。Patel和Kumar\cite{itm2025liveness}提出了结合活体检测的考勤管理方案。Sharma和Gupta\cite{ieee2024raspberry}实现了基于树莓派4B的人脸识别考勤系统。这些研究表明，在端侧设备上运行完整的AI推理管线是可行的，但需要在模型精度和推理速度之间做出权衡。

当前研究趋势体现在以下几个方面。其一，端侧轻量化模型的持续优化，通过模型量化、知识蒸馏等手段降低模型推理的计算开销和内存占用。其二，被动活体检测的深入研究，利用多帧时序信息和微弱生理信号提升检测准确率，减少对用户配合的依赖。其三，多模态融合技术的应用，结合红外、深度、RGB等多种传感器信息提升活体检测的鲁棒性。

\section{研究目标和内容}

针对现有考勤系统在端侧部署、活体检测和识别准确性方面存在的不足，本文在树莓派设备上设计并实现了一套完整的人脸识别考勤系统。本文的主要工作内容包括以下几个方面：

（1）人脸注册流的设计与实现。完成从本地名单导入、用户首次注册改密、正式登录、动作挑战人脸采集到质量校验入库的完整注册链路。设计了基于OpenCV landmarks的姿态估计方法，用于判断抓拍图像是否满足正脸或轻微侧转的动作要求。实现了同人一致性校验，确保挑战过程中始终为同一人。设计了纸张和屏幕翻拍检测，降低注册照被伪造的风险。

（2）实时识别考勤流的设计与实现。完成从摄像头实时采集、人脸检测、活体判断、特征提取、身份比对到签到记录写入的端到端识别链路。设计了严格识别引擎，包含单人检测、人脸面积、编码距离、多帧稳定、时间窗和冷却时间六道门禁。引入OpenVINO anti-spoof-mn3模型进行被动活体前置判断，设计了2秒滑动窗口多帧融合策略，对real\_score进行平滑处理，并引入屏幕重放风险信号分析。

（3）识别后端的优化。采用OpenCV YuNet检测加SFace特征提取的识别链路，通过特征金字塔结构实现多尺度人脸检测，通过仿射变换和余弦相似度匹配实现高精度身份识别。在树莓派CPU上实现了实时推理，单帧处理时间控制在200毫秒以内。

（4）用户门户与远程管理接口的设计与实现。完成基于Jinja2模板的外网用户门户页面，包括注册、登录、用户中心和人脸采集页面。实现设备侧管理员API，提供成员管理、人脸审核、考勤记录查询和设备状态监控接口。设计Windows端远程管理系统通过Tailscale网络访问树莓派API的对接方案。

（5）系统实验与性能分析。对系统的各项功能进行端到端验证，在树莓派实际设备上测试识别帧率、活体推理耗时和内存占用，根据实测结果调整识别阈值和线程策略。对比分析不同配置参数对系统性能的影响。

\section{论文结构}

本文共分为4章，各章主要内容如下。

第1章为绪论，分析了人脸识别考勤系统的研究背景与意义，对国内外人脸识别技术和活体检测方法的研究现状进行了总结与分析，给出了本文的主要工作内容与论文结构安排。

第2章对实现系统的核心技术进行简述，包括YuNet人脸检测模型、SFace特征提取模型、被动活体检测技术、FastAPI框架和SQLite数据库，为后续章节的系统设计与实现奠定技术基础。

第3章进行系统的详细设计与实现，完成五层架构的设计思路阐述，对人脸注册流、实时识别考勤流、识别后端、用户门户与管理接口进行详细讲解，对数据库表结构进行说明。

第4章进行系统实验与性能分析，对注册流程、活体挑战、识别签到和审核机制进行端到端验证，评估系统在树莓派设备上的实际运行性能，分析不同参数配置对识别效果的影响。

最后对全文进行总结，指出系统中尚未完成或需要进一步改进的方面，并对未来工作进行了展望。
